\section{Evaluationsmethode und Ergebnisse}

\subsection{Auswahl der Programmieraufgaben sowie der Referenz- und Fehlerimplementierungen}
Die Programmieraufgaben stammen von der Coding-Challenge-Plattform Leetcode \cite{leetcode_problemset}. Leetcode ist eine öffentliche Website auf der das Lösen verschiedener 
algorithmischer Programmieraufgaben mit unterschiedlichen Verfahren und Datenstrukturen geübt werden kann. Mit der Auswahl der Programmieraufgaben wurde versucht ein breites 
Spektrum an unterschiedlichen Problemarten und möglichen Sonderfällen abzudecken. Es wurden insgesamt sieben Probleme für die Evaluation des Assistentsystems ausgewählt.
Die genauen Aufgabenstellungen und die daraus generierten Testfälle können Anhang~\ref{anhang:programmieraufgaben-testfaelle} entnommen werden.

Die korrekten Referenzimplementierungen sowie die falschen Lösungen zu den aufgeführten Programmieraufgaben stammen aus eigener Entwicklung. Die Referenzimplementierungen 
wurden auf der Leetcode-Plattform eingereicht und gegenüber der dort vorhandenen Testsuite validiert. Dadurch liegt für jede Programmieraufgabe eine geprüft korrekte Lösung vor, 
die im weiteren Verlauf dieser Untersuchung als Vergleichsgrundlage für die Validierung der KI-generierten erwarteten Ergebnisse dient. 

Die hier als bewusst fehlerhaft verwendeten Lösungen basieren ebenfalls auf eigenen früheren Lösungsversuchen zu den jeweiligen Problemen. Für die Untersuchung ist das 
vorteilhaft, da es sich nicht um künstlich konstruierte oder absichtlich eingebaute Fehler, sondern um tatsächilche, während des menschlichen Problemlösungsprozesses 
enstandene Implementierungsfehler handelt. Diese falschen Lösungsversuche liefern für viele typische Testfälle das korrekte Ergebnis, scheitern jedoch an spezifischen Grenz- 
oder Sonderfällen. Dadurch eignen sie sich besonders dafür zu überpfrüfen wie gut die KI-generierten Testfälle fehlerhafte Implementierungen entlarven.

\subsection{Validität der generierten Testfälle}
Ein generierter Testfall kann aus zwei verschiedenen Gründen nicht valide sein. Seine Eingabewerte können außerhalb der in der
Programmieraufgabe definierten Constraints liegen oder der vom Sprachmodell
erzeugte erwartete Sollwert kann falsch sein. In beiden Fällen ist ein Fehlverhalten einer 
Implementierung nicht interpretierbar, wodurch der Testfall unbrauchbar wird. Die Kennzahlen wurden wie folgt berechnet:

\[
Constraint-Konformitaetsrate =
\frac{N_{\mathrm{zulaessig}}}{N_{\mathrm{gesamt}}} \cdot 100
\]

\[
Sollwertvaliditaet =
\frac{N_{\mathrm{zulaessig-und-korrekt}}}{N_{\mathrm{zulaessig}}} \cdot 100
\]

Dabei bezeichnet \(N_{\mathrm{zulaessig}}\) die Anzahl der
constraint-konformen Testfälle und \linebreak \(N_{\mathrm{zulaessig-und-korrekt}}\) die Menge aller zulässigen
Testfälle, deren Sollwert durch die Referenzimplementierung bestätigt wurde.
Ergänzend wurde die Gesamtquote unmittelbar nutzbarer Testfälle berechnet.
Sie bildet das Verhältnis der Testfälle, die gleichzeitig zulässig sind und einen korrekten
Sollwert besitzen, gegenüber allen generierten Testfällen ab.

Zu jede der sieben Programmieraufgabe wurden 15 Testfälle durch das Assistentsystem generiert. 
Insgesamt kommt die Untersuchung somit auf 105 Testfälle.
Zu jeder Aufgabe wurde die korrekte Referenzimplementierung mit den 15 generierten Testfällen ausgeführt.
Das daraus resultierende Ergebnis wurde anschließend mit dem von der KI erwarteten Ergebnis verglichen. 
Die Ergebnisse der Ausführung sowie die dazu genutzten Referenzimplementierungen befinden sich in der Anlage \(test-results.md\).
Das Feld \texttt{matches-}\allowbreak\texttt{expected} im Ergebnisdatensatz gibt an, ob beide Ergebnisse übereinstimmen.

\subsubsection{Ergebnisse der Validitätsprüfung}
Bei der manuelle Constraint-Prüfung konnten insgesamt zehn unzulässige
Testfälle identifiziert werden. Davon entfielen zwei auf Problem 1 und die restlichen 
acht auf Problem 5. Bei Problem 1 wurde mit \(n=2\,863\,311\,530\) ein Wert außerhalb des
erlaubten Bereichs \(1 \leq n \leq 2^{31}-1\) (2\,147\,483\,647) erzeugt. Bei Problem 5
enthielten acht Testfälle mindestens einen Eingabewert 0, obwohl für alle drei
Eingabevariablen positive Werte zwischen 1 und \(10^9\) durch die Beschränkungen 
innerhalb der Aufgabenstellung vorgeschrieben waren.
Damit erfüllten 95 von 105 Eingaben die in den Aufgaben gesetzten Eingabebeschränkungen was 
einer Constraint-Konformitätsrate von 

\[
\frac{95}{105} \cdot 100 = 90{,}48\,\%.
\]

entspricht.

Unter den 95 zulässigen Testfällen wurden 86 Sollwerte durch die
Referenzimplementierungen bestätigt. Neun zulässige Testfälle enthielten
einen falschen Sollwert. Daraus ergibt sich eine Sollwertvalidität von 

\[
\frac{86}{95} \cdot 100 = 90{,}53\,\%.
\]

Bezogen auf alle 105 generierten Testfälle ergibt sich eine
Gesamtquote an unmittelbar nutzbaren Testfällen von

\[
\frac{86}{105} \cdot 100 \approx 81{,}90\,\%.
\]

Die Ergebnisse aus der Auswertung der einzelnen Aufgaben sind in
Tabelle~\ref{tab:validitaet-probleme} zusammengefasst.

\begin{table}[htbp]
    \centering
    \renewcommand{\arraystretch}{1.15}
    \caption{Validität der generierten Testfälle nach Aufgabe}
    \label{tab:validitaet-probleme}
    \begin{tabular}{l r r r r}
        \hline
        Aufgabe & Erzeugt & Zulässig & Korrekt & Sollwertfehler \\
        \hline
        Problem 1: Alternating Bits & 15 & 13 & 13 & 0 \\
        Problem 2: Sort Vowels & 15 & 15 & 7  & 8 \\
        Problem 3: Rotate List & 15 & 15 & 14 & 1 \\
        Problem 4: Highest Altitude & 15 & 15 & 15 & 0 \\
        Problem 5: Minimum Flips & 15 & 7  & 7  & 0 \\
        Problem 6: Sort Colors & 15 & 15 & 15 & 0 \\
        Problem 7: K-Sum Pairs & 15 & 15 & 15 & 0 \\
        \hline
        Gesamt    & 105 & 95 & 86 & 9 \\
        \hline
    \end{tabular}
\end{table}

Die Sollwertfehler konzentrieren sich vollständig auf zwei Aufgaben. Bei
Problem 2, der Sortierung von Vokalen nach Häufigkeit und Erstauftreten,
waren acht der 15 erwarteten Ausgaben falsch. Bei Problem 3, der Rotation einer verketteten Liste, trat ein einzelner
Sollwertfehler in Testfall 12 auf. Die übrigen 14 zulässigen Testfälle
wurden korrekt berechnet. Für alle anderen Probleme erreichten die von der KI generierten Testfälle nach Entfernung
unzulässiger Eingaben eine Sollwertvalidität von \(100\,\%\).

\subsubsection{Interpretation der Untersuchten Validitätsrate}
Aus den Ergebnissen der Validitätsprüfung lassen sich erste interessante Schlüsse ziehen.
Auffällig ist das bei den ausgewählten Aufgaben scheinbar kein Zusammenhang zwischen der 
Constraint-Konformitätsrate und der Sollwertvaliditätsrate besteht. Bei Problem 5 lag die 
Constraint-Konformitätsrate mit \(46{,}66\,\%\) deutlich unter dem Durchschnitt, die Sollwertvaliditätsrate 
der generierten Testfälle zu diesem Problem lag jedoch bei \(100\,\%\). 
Andersrum lag die Constraint-Konformitätsrate bei Problem 2 bei \(100\,\%\) und die Sollwertvaliditätsrate 
nur bei \(46{,}66\,\%\). Die Untersuchung deutet folglich darauf hin, dass eine geringe Constraint-Konformitätsrate 
nicht zwangsläufig eine geringe Sollwertvaliditätsrate impliziert und umgekehrt.

Insbesondere die 8 falsch erzeugten Testfälle für Problem 5 scheinen bezüglich der Constraint-Konformität bemerkenswert. 
Trotz eines klar begrenzten Wertebereiches von 1 bis \(10^9\), enthielten über die Hälfte der generierten Testfälle den Wert 0. 
Dieses Muster könnte darauf zurückzuführen sein, dass der Wert 0 häufig bei der Bildung von Grenz- und Sonderfällen Anwendung findet. 
In diesem Fall wurde das Grenzfallmuster jedoch fehlerhaft auf eine Aufagabe angewandt, deren Aufgabenbeschreibung diesen Fall eindeutig 
und explizit ausschließt.

Die Analyse der Sollwertfehler ergibt eine starke Konzentration der fehlerhaften Erwartungswerte auf Problem 2, 
welche gegen eine rein zufällige Verteilung spricht. 
Die Validität scheint von der komplexität der semantischen Anforderungen der jeweiligen Aufgabe
abzuhängen. Damit ist die Schwierigkeit gemeint, die in der Aufgabenbeschreibung formulierten Regeln korrekt 
zu interpretieren und in ihrer Kombination auf einen konkreten Eingabefall anzuwenden.
Einfache bitweise, kumulative oder partitionierende Operationen wurden in der Stichprobe zuverlässig behandelt, während die
mehrstufige Sortierregel mit positionsabhängiger Tie-Break-Bedingung in Problem 2 deutlich häufiger zu falschen Erwartungen führte. 
Maßgeblich für die Sollwertvalidität scheint somit nicht die die rechnerische Schwierigkeit einer Aufgabe zu sein sondern die Anzahl, 
Abhängigkeit und sprachliche Eindeutigkeit der anzuwendenden Regeln.

Der einzelne Sollwertfehler bei Problem 3 ist insofern interessant, als das bei vergleichbaren Testfällen zu selbiger Aufgabe keine fehlerhaften Werte für das erwartete 
Ergebnis auftraten. 
Aus der Untersuchung lässt sich hierfür keine eindeutige Ursache ableiten. Am ehesten kann der Fehler daher als Veranschaulichung der zuvor beschriebenen Funktionsweise 
von Large Language Models interpretiert werden. Die erzeugte Tokenfolge basiert auf Wahrscheinlichkeitsverteilungen und nicht auf einer garantierten logischen Herleitung. 
Dadurch können auch bei ähnlichen Eingaben vereinzelt fehlerhafte Ausgaben entstehen.


\subsection{Erkennung fehlerhafter Implementierungen}
Nach der Untersuchung der Testfallvalidität wird im zweiten Auswertungsschritt geprüft, ob die
generierten Testfälle tatsächlich dazu geeignet sind, fehlerhafte Implementierungen zu
erkennen. Dafür wurden die generierten Testfälle mit einer gezielt fehlerhaften Lösung ausgeführt. 
Die fehlerhaften Implementierungen zeichnen sich dadurch aus das sie einen vorab bekannten Defekt enthalten, 
welcher jedoch nur bei bestimmten Eingaben zu einem abweichenden Ergebnis führt.
Die Ergenbisse dieser Ausführung sowie die fehlerhaften Implementierungen können ebenfalls der Anlage \(test-results.md\) entnommen werden.
Ein Testfall erkennt eine falsche Implementierungen, wenn das indem Testfall definierte erwartete Ergebnis 
von dem Ergebnis nach Ausführung abweicht. Für diese Analyse werden ausschließlich Testfälle berücksichtigt 
welche zuvor als Valide klassifiziert wurden. 

\subsubsection{Ergebnisse der Fehlererkennung}
Von den 86 validen Testfällen führen 13 bei einer fehlerhaften Implementierung zu einem vom
Sollwert abweichenden Ergebnis. Tabelle~\ref{tab:fehlererkennung-probleme} zeigt, das sich die 
13 Testfälle so auf die einzelnen Aufgaben aufteilen, dass sechs der sieben fehlerhaften Implementierungen 
durch mindestens einen validen Testfall erkannt werden. Dies entspricht einer Erkennungsrate von rund
\(85,71\,\%\). Nur für Problem~5 wurde durch das LLM kein valider Testfall erzeugt, der den eingebauten
Fehler sichtbar macht.

\begin{table}[htbp]
    \centering
    \renewcommand{\arraystretch}{1.15}
    \caption{Erkennung der fehlerhaften Implementierungen durch valide Testfälle}
    \begin{tabular}{p{0.39\textwidth} r r l}
        \hline
        Programmieraufgabe &
        Valide Tests &
        Treffer &
        Fehler erkannt \\
        \hline
        Problem 1: Alternating Bits & 13 & 1 & Ja \\
        Problem 2: Sort Vowels & 7 & 1 & Ja \\
        Problem 3: Rotate List & 14 & 2 & Ja \\
        Problem 4: Highest Altitude & 15 & 5 & Ja \\
        Problem 5: Minimum Flips & 7 & 0 & Nein \\
        Problem 6: Sort Colors & 15 & 3 & Ja \\
        Problem 7: K-Sum Pairs & 15 & 1 & Ja \\
        \hline
        Gesamt & 86 & 13 & 6 von 7 \\
        \hline
    \end{tabular}
    
    \label{tab:fehlererkennung-probleme}
\end{table}

\subsubsection{Einordnung der Ergebnisse}
Bei der Analyse der Ergebnisse zur Fehlererkennung fällt zunächst auf, dass im Durchschnitt bei jeweils 
15 Testfällen pro Programmieraufgabe nur 1,86 Testfälle erzeugt wurden, welche die fehlerhafte Implementierung
erkennen. Dieses Ergebnis scheint auf den ersten Blick niedrig, da im System-Prompt explizit darauf verwiesen wurde, 
insbsondere solche Testfälle mit einzubeziehen. Hierbei darf jedoch nicht außeracht gelassen werden, das die verwendeten 
Fehlerimplementierungen häufig genau durch ein Eingabemuster entlarvt werden. Aus dieser Sichtweise wäre ein Ergebnis von 
einem fehlererkennenden Testfall pro Aufgabe als optimal zu bewerten, da das LLM so nicht weitere Testfälle für den gleichen 
Sonderfall verwendet. Außerdem muss berücksichtigt werden, das nicht ausgeschlossen werden kann, dass unter den verbleibenen 
Testfällen nicht solche dabei sind, die möglicherweise andere, von dem LLM antizipierte Fehlerimplementierungen, identifizieren 
würden.

Die Analyse der Fehlererkennung gepaart mit der zur Validitätsrate impliziert außerdem einen weiteren interessanten Zusammenhang.
Auffällig ist das für Problem 2 ein Testfall generiert wurde welcher die falsche Lösung erkennt, obwohl das LLM zuvor Schwierigkeiten 
damit hatte die korrekten Sollwerte für die eigens generierten Testfälle zu antizipieren. Auf der anderen Seite wurde für Problem 5 kein 
Testfall generiert der die falsche Implementierung erkennen konnte. Hier hatte das LLM vorab Probleme die mitgegebenen Eingabebeschränkungen 
(Constraints) zu befolgen, konnte aber die Sollwerte der Testfälle zu \(100\,\%\) korrekt vorhersagen. Diese Beobachtung schlägt vor, dass 
die Generierung fehlerekennender Testfälle weniger von der semantischen Komplexität der Aufgabenstellung, sondern mehr vom Verständnis der 
vorgegebenen Eingabebeschränkung abhängt.

\newpage